\documentclass[modern]{aastex631}

\usepackage{amsmath}
\usepackage{hyperref}
\usepackage{url}

\begin{document}

\title{GASP: A Corrected, Cross-Matched Catalog of 19,190 
       Gaia DR3 Asteroid Reflectance Spectra}

\author{Werner Scheibenpflug}
\affiliation{Independent researcher, Austria}
\email{office@esoxconsult.com}

\begin{abstract}
We present GASP (Gaia Asteroid Spectral Pipeline), an open-source
pipeline and catalog combining corrected Gaia DR3 asteroid
reflectance spectra with SDSS Moving Object Catalog photometry.
The final catalog contains 19,190 asteroids with 16-band
NUV-corrected spectra (374--1034~nm), SDSS $a^*$ colors,
NEOWISE albedo (22.3\%), family membership (39.3\%, 117 families),
MPC orbital classification (100\%), and spectral slopes $s_1$--$s_4$.
We validate the NUV correction against the Eight Color Asteroid
Survey (ECAS; Zellner et al. 1985), finding Pearson
$r = -0.848$ ($n=83$, $p \approx 0$), confirming consistency
with ground-based photometry. The catalog and pipeline are
publicly available at GitHub and Zenodo
(DOI: \href{https://doi.org/10.5281/zenodo.19366681}
{10.5281/zenodo.19366681}).
\end{abstract}

\keywords{minor planets, asteroids: general --- catalogs ---
techniques: photometric --- surveys}

\section{Introduction}

The third data release of the ESA Gaia mission
\citep[Gaia DR3;][]{GaiaDR3_2023} includes mean reflectance
spectra for 60,518 Solar System objects (primarily asteroids)
across 16 wavelength bands between 374 and 1034~nm. Despite
their scientific potential, these spectra carry a known
systematic artifact: an artificial reddening below 0.55~$\mu$m
caused by inadequate solar analogues in the near-ultraviolet
\citep[NUV;][]{TinautRuano2023}. No publicly available,
corrected, and cross-matched catalog existed prior to this work.

\section{Pipeline}

GASP consists of eight Python scripts executing sequentially.
Starting from the Gaia TAP server \citep{GaiaArchive2024},
we download 968,288 rows (60,518 asteroids $\times$ 16 spectral
bands) and restructure them into wide format (one row per asteroid).
We apply the NUV correction of \citet{TinautRuano2023}, dividing
reflectance values at 374, 418, 462, and 506~nm by factors of
1.07, 1.05, 1.02, and 1.01, respectively (their Table~1).

Cross-matching with the Sloan Digital Sky Survey Moving Object
Catalog \citep[SDSS MOC4;][]{Ivezic2001} on MPC asteroid number
yields 19,190 objects present in both surveys (31.7\% of the
full Gaia SSO catalog). The inclusion of SDSS $a^*$ colors
enables a robust, independent check of C- vs.\ S-complex
separation; combined with the NUV-corrected Gaia slopes, GASP
provides two complementary spectrophotometric discriminators for
asteroid surface composition in a single, consistently processed
dataset.

We further enrich this sample with family membership from the
\citet{Nesvorny2015} HCM catalog (PDS4 archive V2.0, 117
families), spectral taxonomy from \citet{Mahlke2022}, albedo and
diameter from NEOWISE \citep{Masiero2011}, and orbital
classification from the MPC Extended Object Catalog. Spectral
slopes $s_1$--$s_4$ (in $\mu$m$^{-1}$) are computed following
the methodology of \citet{TinautRuano2024}.

\section{Catalog Properties}

Table~\ref{tab:catalog} summarizes the catalog coverage.
Family membership reaches 39.3\% (7,533 objects in 117 families),
with the largest families being Nysa-Polana ($n=3{,}891$),
Vesta ($n=2{,}847$), and Flora ($n=2{,}156$). Spectral taxonomy
is available for 1.9\% of objects (358 asteroids), reflecting
the limited size of currently available ground-based spectral
classifications rather than a pipeline limitation.
NEOWISE albedo and diameter are available for 22.3\% of objects
(4,286 asteroids). All 19,190 objects carry MPC orbital
classification; 94.4\% are Main Belt Asteroids.

\begin{deluxetable}{lc}
\tablecaption{GASP v1 Catalog Summary\label{tab:catalog}}
\tablehead{
  \colhead{Property} & \colhead{Value}
}
\startdata
Total asteroids         & 19,190          \\
Spectral bands          & 16 (374--1034~nm) \\
NUV correction          & Applied (Tinaut-Ruano et al.\ 2023) \\
SDSS photometry         & $u, g, r, i, z$, $a^*$ \\
Family coverage         & 39.3\% (117 families) \\
Taxonomy coverage       & 1.9\% (Mahlke et al.\ 2022) \\
Albedo/diameter         & 22.3\% (NEOWISE) \\
Orbital classification  & 100\% (MPC) \\
ECAS validation         & $r = -0.848$ ($n=83$, $p\approx0$) \\
Catalog columns         & 72 \\
File size               & $\sim$4.9~MB (Apache Parquet) \\
DOI                     & 10.5281/zenodo.19366681 \\
\enddata
\end{deluxetable}

\section{Validation}

We validate the NUV correction by cross-matching our catalog
with the Eight Color Asteroid Survey
\citep[ECAS;][]{Zellner1985}, which provides ground-based
photometry in eight broadband filters including two below
0.45~$\mu$m. Of 589 ECAS asteroids, 83 appear in our catalog.
We compute the Pearson correlation between the ECAS $b-v$ color
index and the GASP reflectance at 418~nm (\texttt{refl\_418}),
the closest Gaia band to the ECAS $b$ filter (central wavelength
0.437~$\mu$m). We find $r = -0.848$ ($p \approx 0$, $n = 83$;
Figure~\ref{fig:ecas}), confirming that the NUV-corrected Gaia
spectra are consistent with independent ground-based photometry.
The negative sign reflects the inverse relationship between
reflectance ratio and color index conventions.

Figure~\ref{fig:scatter} illustrates the scientific content of
the catalog: the SDSS $a^*$ color index
\citep[a C/S-complex separator;][]{Ivezic2002} plotted against
the NUV slope $s_1$, colored by NEOWISE albedo.
C-complex asteroids ($a^* < -0.10$) show systematically lower
albedo (mean $p_V \approx 0.06$) and more negative $s_1$ values,
consistent with UV absorption in hydrated primitive bodies.
S-complex asteroids ($a^* > 0.10$) show higher albedo
(mean $p_V \approx 0.22$) and flatter NUV slopes, consistent
with anhydrous silicate compositions.

\begin{figure}
\plotone{figures/05_ecas_validation.png}
\caption{Left: ECAS $b-v$ color vs.\ GASP \texttt{refl\_418}
(NUV-corrected reflectance at 418~nm) for the 83 asteroids
common to both catalogs. The strong negative correlation
($r = -0.848$, $p \approx 0$) confirms that the NUV correction
produces spectra consistent with ground-based photometry.
Right: Mean spectrum of the ECAS-matched subset compared to the
full GASP sample (gray), with the NUV-corrected region
highlighted (purple shading).\label{fig:ecas}}
\end{figure}

\begin{figure}
\plotone{figures/03_astar_vs_s1.png}
\caption{SDSS $a^*$ color index vs.\ NUV spectral slope $s_1$
(374--418~nm, in $\mu$m$^{-1}$) for 11,052 GASP asteroids with
valid $s_1$ values, colored by NEOWISE geometric albedo $p_V$
where available (gray points: no albedo data). Dashed vertical
lines mark the C/ambiguous ($a^* = -0.10$) and
ambiguous/S ($a^* = +0.10$) boundaries of
\citet{Ivezic2002}. C-complex asteroids (left) show lower
albedo and more negative $s_1$; S-complex asteroids (right)
show higher albedo and flatter NUV
slopes.\label{fig:scatter}}
\end{figure}

\section{Limitations}

Taxonomy coverage reflects the current state of ground-based
spectral surveys, not a pipeline limitation. The SDSS MOC4
cross-match yields 31.7\% of the full Gaia SSO catalog;
objects without SDSS counterparts are available in the
intermediate pipeline output at the same Zenodo record.
A systematic reddening between 0.7 and 0.9~$\mu$m has been
noted in Gaia DR3 SSO spectra \citep{Galinier2023}; a
corresponding NIR correction is not yet implemented and is
planned for a future release with Gaia DR4 data
(expected end of 2026).

\section{Data Availability}

The complete catalog (\texttt{gasp\_catalog\_v1.parquet},
4.9~MB, Apache Parquet format), all pipeline scripts, and
an exploration Jupyter notebook are available at:

\begin{itemize}
\item GitHub: \url{https://github.com/esoxconsult/gasp}
      (MIT license)
\item Zenodo: \url{https://doi.org/10.5281/zenodo.19366681}
\end{itemize}

\begin{acknowledgments}
This research made use of data from the European Space Agency
(ESA) mission Gaia (\url{https://www.cosmos.esa.int/gaia}),
processed by the Gaia Data Processing and Analysis Consortium
(DPAC). This research also made use of the SDSS Moving Object
Catalog \citep{Ivezic2001}, NEOWISE data \citep{Masiero2011},
the Nesvorny family catalog \citep{Nesvorny2015}, and the
Mahlke et al.\ taxonomy \citep{Mahlke2022}. The author thanks
the developers of \texttt{astroquery} \citep{Ginsburg2019},
\texttt{pandas}, \texttt{numpy}, and \texttt{matplotlib}.
\end{acknowledgments}

\begin{thebibliography}{}

\bibitem[Gaia Collaboration et al.(2023)]{GaiaDR3_2023}
Gaia Collaboration, Vallenari, A., Brown, A.~G.~A., et al.\ 2023,
\aap, 674, A1.
\doi{10.1051/0004-6361/202243940}

\bibitem[Gaia Archive(2024)]{GaiaArchive2024}
ESA Gaia Science Archive. \url{https://gea.esac.esa.int/archive/}

\bibitem[Galinier et al.(2023)]{Galinier2023}
Galinier, M., Delbo, M., Avdellidou, C., Galluccio, L.,
\& Marrocchi, Y.\ 2023, \aap, 671, A40.
\doi{10.1051/0004-6361/202245311}

\bibitem[Ginsburg et al.(2019)]{Ginsburg2019}
Ginsburg, A., Sip{\H o}cz, B.~M., Brasseur, C.~E., et al.\ 2019,
\aj, 157, 98.
\doi{10.3847/1538-3881/aafc33}

\bibitem[Ivezi{\'c} et al.(2001)]{Ivezic2001}
Ivezi{\'c}, {\v Z}., Tabachnik, S., Rafikov, R., et al.\ 2001,
\aj, 122, 2749.
\doi{10.1086/323452}

\bibitem[Ivezi{\'c} et al.(2002)]{Ivezic2002}
Ivezi{\'c}, Z., Juri{\'c}, M., Lupton, R.~H., Tabachnik, S.,
\& Quinn, T.\ 2002, Proc.\ SPIE, 4836, 98.
\doi{10.1117/12.460800}

\bibitem[Mahlke et al.(2022)]{Mahlke2022}
Mahlke, M., Carry, B., \& Mattei, P.-A.\ 2022,
\aap, 665, A26.
\doi{10.1051/0004-6361/202243587}

\bibitem[Masiero et al.(2011)]{Masiero2011}
Masiero, J.~R., Mainzer, A.~K., Grav, T., et al.\ 2011,
\apj, 741, 68.
\doi{10.1088/0004-637X/741/2/68}

\bibitem[Nesvorny et al.(2015)]{Nesvorny2015}
Nesvorny, D., Broz, M., \& Carruba, V.\ 2015,
in Asteroids IV, eds.\ P.\ Michel, F.~E.\ DeMeo,
\& W.~F.\ Bottke (Univ.\ Arizona Press), 297.
\doi{10.1088/0004-6256/150/2/48}

\bibitem[Tinaut-Ruano et al.(2023)]{TinautRuano2023}
Tinaut-Ruano, F., Tatsumi, E., Tanga, P., et al.\ 2023,
\aap, 669, L14.
\doi{10.1051/0004-6361/202245134}

\bibitem[Tinaut-Ruano et al.(2024)]{TinautRuano2024}
Tinaut-Ruano, F., Tatsumi, E., de Le{\'o}n, J., et al.\ 2024,
\aap, 686, A76.
\doi{10.1051/0004-6361/202348752}

\bibitem[Zellner et al.(1985)]{Zellner1985}
Zellner, B., Tholen, D.~J., \& Tedesco, E.~F.\ 1985,
\icarus, 61, 355.
\doi{10.1016/0019-1035(85)90133-2}

\end{thebibliography}

\end{document}
